%---------------------------------------------------------------------------------------------------
% CV compatible ATS — Axel Roudaut (FR)
%
% Regles d'analyse automatique (Applicant Tracking Systems) :
%   * Une seule colonne, pas de tableau/tikz/colorbox pour le contenu, pas de photo/icone.
%   * Coordonnees dans le CORPS, jamais en en-tete/pied de page.
%   * Titres standard ; vraies listes a puces ; texte selectionnable ; police standard.
%   * La couleur ne sert qu'a la mise en valeur — tout le contenu reste du vrai texte.
%---------------------------------------------------------------------------------------------------

\documentclass[10pt,a4paper]{article}

\usepackage[utf8]{inputenc}
\usepackage[T1]{fontenc}
\usepackage[french]{babel}
\usepackage{lmodern}
\renewcommand{\familydefault}{\sfdefault}
\usepackage{microtype}

\usepackage[a4paper,top=0.9cm,bottom=0.6cm,left=1.5cm,right=1.5cm]{geometry}
\usepackage{enumitem}
\usepackage{titlesec}
\usepackage{xcolor}
\usepackage[hidelinks]{hyperref}

\definecolor{accent}{HTML}{1F5C8B}
\definecolor{ink}{HTML}{222222}
\definecolor{muted}{HTML}{5A5A5A}
\color{ink}

\setlength{\parindent}{0pt}
\pagestyle{empty}

\setlist[itemize]{leftmargin=1.2em,topsep=0.5pt,itemsep=0.8pt,parsep=0pt,
                  label=\textcolor{accent}{\small$\bullet$}}

\titleformat{\section}
  {\large\bfseries\color{accent}}{}{0pt}{\MakeUppercase}[{\color{accent}\titlerule[0.8pt]}]
\titlespacing{\section}{0pt}{5pt}{3pt}

% Une entree d'experience / formation.
% #1 intitule  #2 organisation  #3 lieu  #4 dates
\newcommand{\cventry}[4]{%
  \textbf{#1}\hfill\textcolor{muted}{\small #4}\\[0.5pt]%
  \textcolor{accent}{\itshape #2}%
  \ifx\relax#3\relax\else\textcolor{muted}{\itshape\,--\,#3}\fi\par\vspace{1.5pt}%
}

\newcommand{\skillrow}[2]{\textbf{#1\,:}\hspace{0.35em}#2\\}

\begin{document}

%---------------------------------------------------------------------------------------------------
% EN-TETE (texte brut, dans le corps)
%---------------------------------------------------------------------------------------------------
{\Huge\bfseries\color{accent}Axel Roudaut}\hspace{0.7em}{\large\color{ink}Ingénieur DevSecOps \& Data Engineer — Cybersécurité}\\[3pt]
{\small\color{muted}%
Montpellier, France \enspace|\enspace
\href{mailto:axel.roudaut@outlook.com}{axel.roudaut@outlook.com} \enspace|\enspace
+33 6 29 76 33 91 \enspace|\enspace
\href{https://www.linkedin.com/in/axel-roudaut-devops/}{linkedin.com/in/axel-roudaut-devops} \enspace|\enspace
\href{https://axelroudaut.github.io/Resume/}{axelroudaut.github.io/Resume}}

%---------------------------------------------------------------------------------------------------
\section{Résumé}
Ingénieur DevSecOps et Data Engineer spécialisé en cybersécurité, plus de 7 ans d'expérience en
conception et exploitation d'infrastructures distribuées et sécurisées (Thales, Ministère des Armées).
Expertise en automatisation CI/CD, orchestration de services et plateformes de traitement de données.
Ouvert à des postes DevOps, Platform Engineering et Data Engineering à Montpellier dès septembre 2026.

%---------------------------------------------------------------------------------------------------
\section{Compétences techniques}
\skillrow{Conteneurs \& orchestration}{Docker, Kubernetes, Helm}
\skillrow{CI/CD \& automatisation}{GitLab CI, Jenkins, Ansible, Terraform, Packer}
\skillrow{Infrastructure \& cloud}{Linux, Windows, VMware, OpenStack, GCP}
\skillrow{Supervision \& observabilité}{Prometheus, Grafana, gestion des logs}
\skillrow{Data \& streaming}{Kafka, Spark, Elasticsearch, ksqlDB, HDFS, Jupyter}
\skillrow{Stockage \& bases de données}{PostgreSQL, MinIO, HDFS}
\skillrow{Sécurité}{SIEM, Suricata, OSSEC, Sysmon, SIGMA, TheHive, auditd}
\skillrow{Développement}{Python, Go, Bash, PowerShell, API REST}

%---------------------------------------------------------------------------------------------------
\vspace{-7pt}
\section{Expériences}

\cventry{Ingénieur DevSecOps Cybersécurité}{Thales}{Rennes, France}{02/2022 -- présent}
\begin{itemize}
  \item Conception et administration de plateformes d'infrastructure basées sur Kubernetes (Kubernetes, Helm, Docker, Ansible, Python).
  \item Déploiement et maintien d'environnements Linux/Windows de simulation d'attaques (VMware API, Packer, Ansible, Terraform, PowerShell).
  \item Exploitation d'une plateforme d'analyse multi-antivirus auto-hébergée, type VirusTotal (Python, Go, Ansible, VMware ESXi, API REST).
  \item Automatisation des pipelines CI/CD pour optimiser les chaînes de production (GitLab CI, Docker, Python, Go, Bash, PowerShell).
  \item Mise en place de la supervision et de l'observabilité des infrastructures (Prometheus, Grafana, Docker, Ansible).
  \item Développement d'un outil de forensic statistique sur des binaires (Python, Jupyter, MinIO, PostgreSQL).
\end{itemize}

\cventry{DevSecOps \& Data Engineer Cyberdéfense}{Ministère des Armées}{Paris, France}{02/2019 -- 02/2022}
\begin{itemize}
  \item Conception et configuration d'une architecture de supervision de sécurité (SIEM) (Suricata, auditd, OSSEC, Sysmon, SIGMA).
  \item Implémentation de pipelines de traitement de données en temps réel (Elasticsearch, Kafka, ksqlDB, Apache Spark, HDFS).
  \item Automatisation des déploiements et de la supervision sur environnements distribués (Docker, Ansible, GitLab, Prometheus, Grafana).
\end{itemize}

\cventry{Stage -- Audit de cybersécurité matérielle}{Trusted Labs}{Meudon, France}{03/2018 -- 09/2018}
\begin{itemize}
  \item Attaques par canaux cachés électroniques pour des pentests de systèmes embarqués (Python, oscilloscope, électronique).
  \item Injection de faute sur noyau Android via manipulation d'horloge (attaque CLKSCREW, Bash, Android).
\end{itemize}

\cventry{Stage -- Ingénieur d'intégration cloud}{Altair}{Antony, France}{06/2017 -- 08/2017}
\begin{itemize}
  \item Développement d'un connecteur Python d'intégration cloud pour applications HPC (Python, CI/CD, Jenkins, GCP).
\end{itemize}

%---------------------------------------------------------------------------------------------------
\section{Formation}

\cventry{Mastère spécialisé (M2) -- Sécurité des systèmes d'information et des réseaux}{TLS-SEC, Toulouse}{}{2018}
\vspace{-2pt}{\small\color{muted}Cryptographie, assembleur, reverse engineering, malware, noyau Linux, dépassements de tampon, sécurité réseau.}\par\vspace{2pt}

\cventry{Diplôme d'ingénieur -- Télécommunications \& Réseaux}{ENSEEIHT, Toulouse}{}{2015 -- 2018}
\vspace{-2pt}{\small\color{muted}Java, C, assembleur, protocoles réseau, Cisco, algorithmes de routage, transmission du signal.}\par\vspace{2pt}

\cventry{Échange Erasmus -- Big Data \& cybersécurité}{Dublin City University (DCU), Irlande}{}{2017}
\vspace{-2pt}{\small\color{muted}Cryptographie, forensic, sécurité réseau, reverse engineering, Big Data (Hadoop).}\par\vspace{2pt}

\cventry{Classes préparatoires (CPGE) -- Physique \& Sciences de l'Ingénieur}{Lycée Michelet, Vanves}{}{2013 -- 2015}

%---------------------------------------------------------------------------------------------------
\vspace{4pt}
{\color{accent}\titlerule[0.8pt]}\vspace{4pt}
\textbf{\color{accent}Langues :}\enspace Français (langue maternelle), Anglais (professionnel)%
\hfill
\textbf{\color{accent}Permis :}\enspace Moto (A), Voiture (B)

\end{document}
